% A0 cheat sheet – your team's SWEB code (osw26e3 main, team/BASE)
% build: pdflatex cheatsheet.tex  (twice not needed)
\documentclass[10pt,a4paper,landscape]{article}
\usepackage[T1]{fontenc}
\usepackage[utf8]{inputenc}
\usepackage[landscape,margin=0.9cm,top=0.6cm,bottom=0.7cm]{geometry}
\usepackage{lmodern}
\usepackage{xcolor}
\usepackage{tcolorbox}
\tcbuselibrary{skins,breakable}
\usepackage{multicol}
\usepackage{tabularx}
\usepackage{array}
\usepackage{listings}
\usepackage{enumitem}
\usepackage{amssymb}
\usepackage[hidelinks]{hyperref}

\pagestyle{empty}
\setlength{\parindent}{0pt}
\setlength{\columnsep}{0.6cm}
\setlist{nosep,leftmargin=1.1em,itemsep=1pt}
\renewcommand{\arraystretch}{1.15}

% ---------- colours ----------
\definecolor{cBlue}{HTML}{2563EB}
\definecolor{cGreen}{HTML}{16A34A}
\definecolor{cOrange}{HTML}{EA580C}
\definecolor{cPurple}{HTML}{7C3AED}
\definecolor{cTeal}{HTML}{0D9488}
\definecolor{cPink}{HTML}{DB2777}
\definecolor{cRed}{HTML}{DC2626}
\definecolor{cSlate}{HTML}{475569}
\definecolor{cAmber}{HTML}{B45309}
\definecolor{codeBg}{HTML}{F8FAFC}
\definecolor{kw}{HTML}{1D4ED8}
\definecolor{cmt}{HTML}{6B7280}
\definecolor{str}{HTML}{B45309}

% ---------- code ----------
\lstdefinestyle{sweb}{
  language=C++, basicstyle=\ttfamily\scriptsize, columns=fullflexible, keepspaces=true,
  keywordstyle=\color{kw}, commentstyle=\color{cmt}\itshape, stringstyle=\color{str},
  morekeywords={size_t,uint64,nodiscard,ScopeLock,Mutex,UserProcess,UserThread,ArchMemory,
    ArchMemoryMapping,PageManager,Loader},
  backgroundcolor=\color{codeBg}, frame=none, xleftmargin=2pt, aboveskip=2pt, belowskip=2pt,
  showstringspaces=false, literate={~}{{\textasciitilde}}1
}
\lstdefinestyle{shell}{
  basicstyle=\ttfamily\scriptsize, columns=fullflexible, keepspaces=true,
  backgroundcolor=\color{codeBg}, commentstyle=\color{cmt}\itshape, morecomment=[l]{\#},
  aboveskip=2pt, belowskip=2pt, showstringspaces=false
}
\lstset{style=sweb}
\newcommand{\cc}[1]{\texttt{#1}}   % inline code
\newcommand{\ok}{\textcolor{cGreen}{\checkmark}}
\newcommand{\no}{\textcolor{cRed}{\ensuremath{\times}}}

% ---------- boxes ----------
\newtcolorbox{csec}[2]{enhanced, breakable, colback=#1!5, colframe=#1, coltitle=white,
  fonttitle=\bfseries\small, title={#2}, boxrule=0.6pt, arc=2.5mm,
  left=4pt, right=4pt, top=3pt, bottom=3pt, toptitle=1.5pt, bottomtitle=1.5pt,
  before skip=3pt, after skip=3pt, fontupper=\small\raggedright}
\newtcolorbox{tip}[1][cAmber]{enhanced, colback=#1!10, colframe=#1!70, boxrule=0.4pt,
  arc=1.5mm, left=3pt, right=3pt, top=1pt, bottom=1pt, before skip=3pt, after skip=3pt,
  fontupper=\footnotesize}

\newcolumntype{Y}{>{\raggedright\arraybackslash}X}
\newenvironment{tab}[1]{\footnotesize\tabularx{\linewidth}{#1}}{\endtabularx}

\begin{document}

% ================= header =================
\begin{tcolorbox}[enhanced, colback=cBlue, colframe=cBlue, arc=3mm, boxrule=0pt,
  interior style={left color=cBlue, right color=cPurple}, left=8pt, top=5pt, bottom=5pt]
  \color{white}{\LARGE\bfseries A0 Cheat Sheet}\hfill{\small team osw26e3 \textbullet{} SWEB x86-64}\\[2pt]
  {\small Your team's code (\texttt{main} of 13 Oct, \texttt{40140475}) \textbullet{} the API you need for the
  interview \textbullet{} read it once, keep it next to CLion}
\end{tcolorbox}

\begin{multicols}{2}

% ================= 1 access =================
\begin{csec}{cBlue}{1 \quad How to get to everything}
\begin{lstlisting}
UserProcess* process = currentThread->getProcess();
ArchMemory& mem   = process->getAddressSpace(); // page tables
Mutex& lock       = process->getThreadLock();   // per-process lock
Loader& loader    = process->getLoader();       // ELF loader
UserThread* ut    = (UserThread*) currentThread; // UserThread fields
\end{lstlisting}
\begin{itemize}
  \item \cc{currentThread} = the thread running right now (global \cc{Thread*}).
  \item Pointer $\to$ \cc{->}, reference $\to$ \cc{.}:\\
        \cc{process->getAddressSpace().mapPage(...)}
\end{itemize}
\begin{tip}[cRed]
  \no{} \cc{currentThread->loader\_}, \cc{(UserProcess*) currentThread} = the tutorials' (upstream) style - \textbf{not} your team's code.
\end{tip}
\end{csec}

% ================= 2 constants & formulas =================
\begin{csec}{cGreen}{2 \quad Pages: constants \& formulas}
\begin{tab}{@{}l l Y@{}}
\cc{PAGE\_SIZE}      & 4096 & one page = one frame \\
\cc{USER\_BREAK}     & \cc{0x800000000000ULL} & below = user, above = kernel \\
\cc{USER\_BREAK\_VPN} & \cc{USER\_BREAK/PAGE\_SIZE} & same, as page number \\
\end{tab}
\begin{lstlisting}
size_t vpn    = address / PAGE_SIZE;   // address -> page number
size_t addr   = vpn * PAGE_SIZE;       // page number -> address
size_t offset = address % PAGE_SIZE;   // position inside the page
bool aligned  = (address % PAGE_SIZE == 0);
// next page: vpn + 1   (NOT vpn + PAGE_SIZE!)
\end{lstlisting}
\begin{tip}[cGreen]
  \textbf{vpn} = virtual page (what the program sees) \quad
  \textbf{ppn} = physical frame (real RAM)\\
  Tree: PML4 $\to$ PDPT $\to$ PD $\to$ PT, 512 entries each. TLB = cache of translations.
\end{tip}
\end{csec}

% ================= 3 ArchMemory =================
\begin{csec}{cTeal}{3 \quad ArchMemory \quad\normalfont(= \cc{process->getAddressSpace()})}
\begin{tab}{@{}>{\raggedright}p{0.43\linewidth} Y@{}}
\cc{mapPage(vpn, ppn, 1)} & vpn $\to$ ppn, user access. \textbf{false} if already mapped. \cc{[[nodiscard]]}! \\
\cc{unmapPage(vpn)} & removes mapping, \textbf{frees the ppn itself}, no TLB flush \\
\cc{resolveMapping(vpn)} & $\to$ \cc{ArchMemoryMapping} (what is mapped?) \\
\cc{checkAddressValid(addr)} & $\neq 0$ if mapped \\
\cc{ArchMemory::flushTlb()} & \textit{static} -- after every \cc{unmapPage} \\
\cc{ArchMemory::getIdent\-AddressOfPPN(ppn)} & \textit{static} -- kernel pointer into a frame \\
\end{tab}

\medskip
\textbf{ArchMemoryMapping} -- \textit{fields, no} \cc{()}
\begin{tab}{@{}l Y@{}}
\cc{.page\_size} & \textbf{0 = not mapped}, else 4096 \\
\cc{.page\_ppn}  & the frame it points to \\
\cc{.pt\_ppn .pd\_ppn .pdpt\_ppn .pml4\_ppn} & page-table pages on the way \\
\end{tab}
\begin{lstlisting}
ArchMemoryMapping m = process->getAddressSpace().resolveMapping(vpn);
if (m.page_size == 0) { /* free slot */ }
\end{lstlisting}
\end{csec}

% ================= 4 PageManager =================
\begin{csec}{cGreen}{4 \quad PageManager \quad\normalfont(physical frames)}
\begin{lstlisting}
size_t ppn = PageManager::instance()->allocPPN(); // free + ZEROED
PageManager::instance()->freePPN(ppn);    // only if NOT mapped
PageManager::instance()->getNumFreePages();
PageManager::instance()->getTotalNumPages();
\end{lstlisting}
\begin{tip}[cGreen]
  \cc{allocPPN} = get a frame, \cc{mapPage} = connect it. You need \textbf{both}, in this order.
  \cc{allocPPN()} takes \textbf{no} vpn (its parameter is a page \emph{size}).
\end{tip}
\end{csec}

% ================= 5 locks =================
\begin{csec}{cOrange}{5 \quad Locks}
\begin{lstlisting}
{
  ScopeLock lock(process->getThreadLock()); // acquire ...
  // check + change in ONE critical section
}                                           // ... released at }
\end{lstlisting}
\begin{itemize}
  \item \cc{ScopeLock} releases on \textbf{every} \cc{return} automatically.
  \item By hand: \cc{acquire()} / \cc{release()} -- then release before \textbf{every} \cc{return}.
  \item \cc{Condition}: \cc{wait()}, \cc{signal()}, \cc{broadcast()} -- mutex must be held.
  \item \cc{allocPPN} \emph{before} the lock = shorter critical section
        (give back with \cc{freePPN} on failure).
  \item Never allocate while holding a \cc{SpinLock}.
\end{itemize}
\begin{tip}[cOrange]
  \cc{thread\_lock\_} protects threads, stack slots \& page tables. The page fault handler
  holds it for stack pages but calls \cc{loadPage} \textbf{without} it.
\end{tip}
\end{csec}

% ================= 6 threads =================
\begin{csec}{cPurple}{6 \quad Threads \& registers}
\begin{tab}{@{}>{\raggedright}p{0.5\linewidth} Y@{}}
\cc{currentThread->getTID()} & thread ID \\
\cc{process->getPID()}, \cc{getPPID()} & process IDs \\
\cc{currentThread->user\_registers\_->rip} & where user code continues \\
\cc{->rsp} & user stack pointer \\
\cc{->rdi}, \cc{->rsi}, \cc{->rdx} & 1st, 2nd, 3rd argument \\
\cc{->rax} & return value \\
\cc{process->createThread(entry,false,0,0)} & new \cc{UserThread} \\
\cc{process->findThreadDataByTid(tid)} & \cc{UserThreadData*} \\
\cc{process->exitProcess()} & kills \textbf{all} threads (POSIX) \\
\end{tab}

\smallskip
\cc{UserThreadData}: \cc{thread\_}, \cc{exit\_value\_}, \cc{tid\_}, \cc{exited\_}, \cc{joined\_}, \cc{requested\_cancel\_}\\
\cc{UserThread}: \cc{user\_stack\_top\_}, \cc{lowest\_mapped\_stack\_vpn\_}
\end{csec}

% ================= 7 loader / PFH =================
\begin{csec}{cTeal}{7 \quad Loader \& page faults}
\begin{itemize}
  \item \cc{common/source/mm/PageFaultHandler.cpp}: valid branch $\to$ stack growth
        (\cc{isUserStackFault}, under the lock) $\to$\\
        \cc{getLoader().loadPage(address, getAddressSpace())} (no lock).
  \item New special regions (heap, magic page) go \textbf{first} in the valid branch, then \cc{return}.
  \item \textbf{ELF}: sections (\cc{.text .rodata .data .bss}) live in LOAD segments;
        \cc{.bss} = MemSiz $>$ FileSiz (zeroed).
\end{itemize}
\begin{tab}{@{}l Y@{}}
exit \textbf{666}  & no ELF segment for the address (\cc{loadPage}) \\
exit \textbf{9999} & invalid fault (null pointer, kernel address) \\
\end{tab}
\end{csec}

% ================= 8 six places =================
\begin{csec}{cPink}{8 \quad New syscall: the six places}
\begin{tab}{@{}c >{\raggedright}p{0.36\linewidth} Y@{}}
1 & \cc{syscall-definitions.h} & \cc{\#define sc\_foo 1505} \\
2 & \cc{Syscall.h} & \cc{static size\_t foo(size\_t arg);} \\
3 & \cc{Syscall.cpp} switch & \cc{case sc\_foo: return\_value = foo(arg1); \textbf{break};} \\
4 & \cc{Syscall.cpp} & \cc{size\_t Syscall::foo(size\_t arg)\{...\}} \\
5 & \cc{userspace/libc/include/nonstd.h} & \cc{int foo(void* p);} \\
6 & \cc{userspace/libc/src/nonstd.c} & \cc{\_\_syscall(sc\_foo, (size\_t)p, 0,0,0,0)} \\
\end{tab}
\begin{tip}[cPink]
  Kernel headers: \cc{common/include/kernel/}. Your block: \textbf{1500--1999}
  (taken: 1500 mkdir, 1501 join, 1502 cancel, 1503 testcancel, 1998/1999).
  Also taken: 1000--1002 pthread\_create/exit/self, 2000--2002 fork/getpid/getppid, 2500 execv.
  Free: 1505, 1510/1511, \dots{} Then a test \cc{userspace/tests/foo\_basic.c}.
\end{tip}
\end{csec}

% ================= 9 user pointers =================
\begin{csec}{cRed}{9 \quad User pointers \& return values}
\begin{lstlisting}
// ALWAYS before the kernel touches user memory (no overflow!):
if (buffer == 0 || size > USER_BREAK || buffer > USER_BREAK - size)
  return (size_t) -1;
Syscall::copyUserString(user_ptr, kernel_buf, buf_size); // strings
\end{lstlisting}
\begin{tab}{@{}l l Y@{}}
error & \cc{return (size\_t) -1;} & \no{} \cc{-1U} = only \cc{0xFFFFFFFF} \\
NULL  & \cc{return 0;} & \no{} \cc{NULL} / \cc{-1} in a \cc{size\_t} function \\
types & \cc{size\_t} & for addresses \& sizes, never \cc{int} \\
\end{tab}
\end{csec}

% ================= 10 debug =================
\begin{csec}{cSlate}{10 \quad Debug output}
\begin{lstlisting}
debug(SYSCALL, "vpn %zx ppn %zx\n", vpn, ppn); // %zx hex, %zd dec
currentThread->printBacktrace();               // who called whom
// own flag in common/include/console/debug.h:
const size_t NOTE = Ansi_Cyan | OUTPUT_ENABLED;
\end{lstlisting}
Flags: \cc{SYSCALL PAGEFAULT LOADER USERPROCESS USERTHREAD}
\end{csec}

% ================= 11 shell =================
\begin{csec}{cSlate}{11 \quad Shell}
\begin{lstlisting}[style=shell]
task start a0-02   # new task branch (training copy)
task test          # build + run the test
task run           # just run SWEB
task debug         # QEMU waits; CLion: 127.0.0.1:1234,
                   #   symbols /tmp/sweb-ag/kernel64.x
task diff a0-02    # compare with the solution
# interview = team repo, plain SWEB:
mkdir -p /tmp/sweb && cd /tmp/sweb
cmake ~/Documents/University/2026-fall/os/repos/osw26e3
make && make qemu
\end{lstlisting}
\end{csec}

% ================= 12 pattern =================
\begin{csec}{cBlue}{12 \quad Pattern: find a free slot + map it (A0-02)}
\begin{lstlisting}
size_t Syscall::allocPage()
{
  UserProcess* process = currentThread->getProcess();
  ScopeLock lock(process->getThreadLock());
  for (size_t i = 0; i < ALLOC_REGION_PAGES; i++)
  {
    size_t vpn = ALLOC_REGION_START / PAGE_SIZE + i;
    if (process->getAddressSpace().resolveMapping(vpn).page_size == 0)
    {
      size_t ppn = PageManager::instance()->allocPPN();
      bool mapped = process->getAddressSpace().mapPage(vpn, ppn, 1);
      assert(mapped);
      return vpn * PAGE_SIZE;            // the ADDRESS, not the ppn
    }
  }
  return 0;                              // full -> NULL
}
// freePage: check aligned + in region + mapped (else -1),
//   unmapPage(vpn); ArchMemory::flushTlb(); return 0;   no freePPN!
\end{lstlisting}
\end{csec}

% ================= 13 traps =================
\begin{csec}{cRed}{13 \quad Classic traps}
\begin{enumerate}[label=\textbf{\arabic*}, leftmargin=1.3em]
  \item \cc{case} without \cc{break}
  \item \cc{mapPage} result ignored (\cc{[[nodiscard]]})
  \item \cc{freePPN} after \cc{unmapPage} $\to$ double free
  \item no \cc{flushTlb()} after \cc{unmapPage}
  \item check and map in two lock sections $\to$ race
  \item lock still held at a \cc{return} (manual acquire)
  \item \cc{vpn + i * PAGE\_SIZE} instead of \cc{vpn + i}
  \item \cc{-1} / \cc{-1U} / \cc{NULL} instead of \cc{(size\_t)-1} / \cc{0}
  \item \cc{m.page\_size()} -- it's a field: \cc{m.page\_size}
  \item \cc{int} for addresses or sizes
  \item \cc{allocPPN} before checking the slot $\to$ leak
\end{enumerate}
\end{csec}

% ================= 14 tutor =================
\begin{csec}{cPurple}{14 \quad Tutor questions -- say it out loud}
\begin{itemize}
  \item What does your lock protect? Show the critical section.
  \item Who frees the frame / the page-table pages? (\cc{unmapPage})
  \item Why flush the TLB? Would the test notice?
  \item What if \cc{mapPage} returns false?
  \item What happens to pages never freed at exit? (\cc{\textasciitilde UserProcess} deletes the address space)
  \item Why \cc{size\_t} and \cc{(size\_t)-1}? Why the overflow-safe pointer check?
\end{itemize}
\end{csec}

% ================= 15 libc only =================
\begin{csec}{cGreen}{15 \quad Not every task is a syscall (A0-13 \texttt{atexit})}
Needs nothing only the kernel has? $\to$ only the libc stub:
\begin{lstlisting}
// userspace/libc/src/stdlib.c                (exit() is in exec.c)
static void (*handlers[32])(void);  static int count = 0;
int atexit(void (*f)(void))
{ if (f == 0 || count == 32) return -1;  handlers[count++] = f;  return 0; }
void __runAtexitFunctions(void)          // exit() calls it before sc_exit
{ while (count > 0) { count--; handlers[count](); } }   // reverse, pop first
\end{lstlisting}
Missing: \cc{strcpy/strncpy/strcat} (A0-14). Stubs: \cc{malloc}, \cc{sem\_*}, \cc{pthread\_mutex/spin/cond\_*}, \cc{sleep}, \cc{waitpid}, \cc{mmap} (\cc{ATLAS.md} \S 5).
\end{csec}

\end{multicols}
\end{document}
